\documentclass[conference]{IEEEtran}

\usepackage[utf8]{inputenc}
\usepackage[spanish, english]{babel}
\usepackage{amsmath}
\usepackage{amssymb}

\title{Sistema de Reconocimiento Facial basado en Eigenfaces mediante PCA y SVD}

\author{
    \IEEEauthorblockN{Ipanaqué Torres, Jesús Andrés}
    \IEEEauthorblockA{
        Facultad de Ciencias\\
        Universidad Nacional de Ingeniería\\
        Lima, Perú
    }
}


\begin{document}

\maketitle

\begin{abstract}
Los sistemas de reconocimiento facial modernos basados en redes neuronales 
presentan evidentes limitaciones en cuanto a interpretación matemática y costo computacional.
En el presente trabajo, se propone un sistema de reconocimiento facial basado en el algoritmo 
de Eigenfaces, desarrollado por Turk y Pentland (1991), haciendo uso de Principal Component 
Analysis (PCA) a través de la Singular Value Decomposition (SVD).
La implementación se realizará en Python haciendo uso exclusivo de NumPy, evitando librerías 
de alto nivel, con el fin de exponer la estructura algebraica del reconocimiento de patrones faciales.
Se emplea la base de datos AT\&T Database of Faces para la evaluación del sistema.
El sistema representa cada imagen como un vector en un espacio de alta dimensión, 
calcula el ``rostro promedio'' y centra los datos para luego aplicar SVD sobre la 
matriz de rostros centrados, obteniendo las componentes principales denominadas eigenfaces, 
que capturan las direcciones de máxima varianza.
Un ``rostro nuevo'' es reconocido proyectándolo al subespacio reducido y calculando 
su distancia euclidiana respecto a los rostros de la base de datos.
Se evalúa la tasa de reconocimiento en función del número de componentes principales 
retenidas $k$, analizando el equilibrio entre precisión y reducción dimensional.
Se espera demostrar que un sistema matemáticamente interpretable y computacionalmente 
ligero puede alcanzar tasas de reconocimiento aceptables sobre la base de datos.
\end{abstract}

\section{Introduction}

El reconocimiento facial es un problema clásico de la visión computacional,
pues trae consigo aplicaciones en seguridad, autenticación y análisis biométrico.
En la actualidad, el avance de las redes neuronales profundas
ha permitido lograr tasas de reconocimiento altas que antes eran inimaginables;
sin embargo, esto lleva consigo una limitación estructural fundamental:
funcionamiento interno ininterpretable, lo que dificulta en gran medida el análisis, corrección y aplicación del mismo
y más aún en contextos en los cuales es necesaria la transparencia. En sistemas de identificación
judicial, si el sistema de reconocimiento facial es usado para identificar
a una persona como sospechosa, es posible que al querer impugnar alguna decisión, resulta
imposible hacerlo con un sistema basado en redes neuronales, pues
no hay estructura matemática que auditar ni forma de demostrar cómo y por qué el sistema falló
obligando así a depender de pruebas empíricas exhaustivas en lugar de un análisis formal.
En sistemas de infraestructura crítica, como el control a plantas de energía,
instalaciones militares o bancos, es un requisito que sea posible auditarlo de manera formal
y no a base de experimentos. En el ámbito de la investigación, en el cual,
generalmente, para construir sobre un resultado es necesario saber el cómo y
por qué funciona, un modelo ininterpretable no permite derivar mejoras teóricas.

El presente trabajo propone un sistema alternativo basado en álgebra lineal: el algoritmo Eigenfaces
planteado por Turk y Pentland en 1991 como una de las soluciones pioneras computacionalmente viables
al problema del reconocimiento facial teniendo en cuenta la limitada capacidad computacional de ese
entonces al reemplazar la comparación pixel por pixel a la proyección de rostros en un subespacio
con menor dimensión.
A diferencia de los sistemas que hacen uso de las redes neuronales, los cuales se denominan
como ``caja negra'' debido a su ininterpretabilidad, Eigenfaces exhibe los patrones de variación
utilizados para el reconocimiento de rostros, permitiendo así el análisis e interpretación
del proceso de reconocimiento. El sistema es construido en Python haciendo uso exclusivo de la
librería Numpy con el fin de conservar la mayor transparencia posible de la estructura algebraica.
En la actualidad, pese a ser un sistema menos preciso que el sistema basado en redes neuronales,
Eigenfaces conserva importancia en ser aprendido como modelo de estudio, pues al exponer de
manera explícita los fundamentos matemáticos del reconocimiento de patrones (reducción 
dimensional, análisis de varianza, proyección en subespacios, ...) permite construir 
intuición sobre el problema antes de abordar sistemas con mayor ininterpretabilidad, permitiendo
así una comprensión más profunda de los conceptos para luego abordar sistemas de mayor complejidad
y no limitarse a solo aplicar lo ya establecido.

Se establecen los siguientes objetivos específicos para el presente trabajo:
(1) implementación del algoritmo Eigenfaces haciendo uso exclusivo de operaciones
de álgebra lineal sobre Numpy;
(2) análisis de las implicaciones de los números de componentes utilizados $k$
sobre la tasa de reconocimiento del sistema; y
(3) demostración de que un modelo matemáticamente interpretable es viable
frente a sistemas de mayor complejidad computacional (como es el caso de las redes neuronales).
Los objetivos en cuestión buscan demostrar que la interpretabilidad matemática 
y la eficiencia computacional no son excluyentes, que un sistema desarrollado
a base de fundamentos matemáticos sólidos puede constituir una herramienta
válida tanto para su aplicación como su uso para la formación y comprensión de los
principios del reconocimiento facial.

El presente sistema se limita al reconocimiento facial en condiciones controladas
haciendo uso de la base de datos AT\&T Database of Faces, el cual contiene 400 imágenes
de 40 sujetos diferenciables con iluminación y pose uniforme; por tanto, no se abordará
el reconocimiento facial en tiempo real ni la previa detección facial, asimismo, se asume
que las imágenes de entrada ya están recortadas y alineadas.
Además, el sistema se enmarca en el contexto del aprendizaje automático clásico,
en el subcampo de reconocimiento de patrones basado en métodos algebraicos y
estadísticos, específicamente. El sistema trabaja eficientemente con conjuntos de
datos pequeños a diferencia de otros enfoques modernos que, para buenos resultados,
requieren obligatoriamente de grandes volúmenes de datos y recursos computacionales,
pues para diferenciar rostros no siguen un procedimiento matemático específico,
sino que ``aprenden'' a base de prueba y error; por tanto, el sistema planteado
ofrece una alternativa viable en escenarios con recursos limitados y esta ventaja
cobra relevancia en aquellos donde la recopilación masiva de datos biométricos presenta
restricciones éticas o logísticas. Para finalizar, es importante mencionar que 
el presente trabajo no busca competir con otros modelos en materia de precisión,
sino demostrar que es posible construir un sistema de reconocimiento facial viable
a partir de herramientas matemáticas elementales. El valor del trabajo consiste
en la transparencia del proceso al mostrar de manera explícita la justificación
algebraica, hacer posible el análisis de cada resultado y predecible cada limitación
generada por la estructura del sistema.

\end{document}